%!TEX root = ../main.tex
\section{Open Challenges and Research Agenda}
\label{sec:agenda}

The taxonomy identifies gaps that are algorithmically concrete, not generic calls
for ``more data'' or ``better generalization.''  We frame each challenge as an
estimand or system property and state how progress could be falsified.

\subsection{Calibrating Proxies to Policy Utility}

Most preparation runs cannot afford retraining a policy for every candidate
decision.  The central estimation problem is to predict
$V(e\mid c,\mathcal{D})$ from cheap features while preserving ranking under
changes in learner, budget, and target distribution.  Current proxies---smoothness,
information, progress, influence, or foundation-model judgments---capture
different projections and can disagree with downstream outcomes.

A promising design is a hierarchical calibrator that combines technical and
semantic features with a small, actively selected set of retraining or rollout
observations.  It should output uncertainty and decompose value into correctness,
coverage, transfer, and risk rather than one opaque score.  Evaluation should
measure rank calibration and regret of the resulting data decision across held-out
learners and tasks, not only correlation on the calibration set.  An estimator is
useful if it achieves near-oracle fixed-budget policy utility with substantially
fewer retraining/robot queries and degrades gracefully under target shift.

\subsection{Estimating Interaction and Marginal Value Below the Episode}

Individual scores miss complementarity.  A grasp segment may be useless without
its approach; two sources can conflict despite each being useful alone; a rare
failure becomes valuable only when paired with a recovery.  Exact Shapley or
leave-one-out retraining is infeasible at VLA scale, while pointwise influence
approximations can ignore non-convex training and sequence dependence.

Research should model value over hierarchical groups and interactions:
transition $\rightarrow$ primitive $\rightarrow$ episode $\rightarrow$ task
$\rightarrow$ source.  Sparse group influence, low-rank datamodels, gradient
sketches, and counterfactual curriculum experiments are candidate tools.  A
benchmark can plant known complementary and conflicting groups and evaluate
whether selection preserves the combination under equal compute.  This directly
tests whether an algorithm has learned data structure rather than a visible
quality artifact.

\subsection{Learnable Cross-Embodiment Action Semantics}

The field lacks a standard that expresses what an action \emph{does} across
morphologies and controllers.  Fixed action vectors and masks solve serialization,
not consequence alignment.  A stronger representation would condition on a
kinematic graph, end-effector and controller specification, reference frame,
latency, and uncertainty, and map actions through a shared transition-level
latent.

The representation should be tested by cycle consistency and physical outcome:
map a source action to a target embodiment, execute or simulate it, and predict
the same object-centric state change within tolerance.  Held-out morphology and
controller families are essential.  A mapping that improves cross-robot policy
loss but violates target limits or contact should fail the benchmark.  Data formats
should preserve original actions and reversible transformations so a learned
standard can replace an earlier one without recollection.

\subsection{Verification of Invisible Physical Errors}

Generative models are optimized for observable likelihood or preference and can
produce errors invisible in pixels: inconsistent depth across views, contact
without force closure, an unreachable wrist pose, or an action label that does not
cause the shown transition.  The research problem is to estimate the probability
that a generated unit is usable for a stated training role.

Progress requires adversarial test suites with paired plausible/invalid examples
at each validation layer, including subtle contact and controller errors.
Verifiers should be evaluated on calibration under new generators and assets,
not just accuracy on their own generator.  Generator--verifier dependence must be
reported.  Selective promotion, in which uncertain samples are routed to
simulation or real replay, can be measured by verified yield, downstream utility,
and cost.  This turns ``synthetic data quality'' into a testable cascade-design
problem.

\subsection{A Schema for Failure, Recovery, and Safety}

Current datasets use inconsistent episode-level success fields and often omit
takeover, retry, near miss, or failure cause.  A shared schema should record:
outcome definition and confidence; first deviation or failure interval;
precondition violated; human/automatic intervention; retry linkage; valid prefix;
recovery strategy and result; safety severity; and whether actions remain valid
targets.  It must allow multiple causes and uncertain temporal boundaries.

The schema's value can be evaluated by holding raw trajectories fixed and testing
whether structured routing improves recovery and worst-case safety over success-only,
all-data, and binary failure baselines.  Inter-annotator/model agreement and
cross-dataset mapping are necessary but not sufficient; the decisive outcome is
whether a policy learns to identify and exit failure regions without imitating
unsafe actions.

\subsection{Stable, Safe Data Self-Evolution}

A self-evolving engine is a feedback controller with delayed, biased, and costly
observations.  Open questions include: How rapidly may mixture weights change?
Which anchor distribution prevents forgetting?  How can verifier drift be
detected without a fresh oracle?  When should the system collect, generate,
relabel, or stop?  How can safety constraints be enforced while acquiring
information?

Evaluation should run multi-round closed-loop benchmarks, not one post-training
iteration.  Metrics include cumulative robot cost, area under the utility curve,
worst-group retention, calibration drift, error amplification, safety events, and
rollback frequency.  Stress tests should seed a biased verifier or generator and
measure whether the loop detects and corrects the bias.  A method that obtains a
short-term average gain while collapsing rare capabilities is unstable, even if
its final aggregate success is high.  The recursive-data collapse result is not a
robotics theorem, but it supplies an actionable stress test: progressively
increase the fraction of self-generated experience while measuring tail coverage
and keep-versus-replace policies under a fixed real anchor
~\cite{shumailov2024collapse}.  RECAP's checkpoint restart, SOP's mixture bounds,
and LWD's offline reference provide concrete controls to ablate rather than
informal assurances~\cite{intelligence20250_2511_14759,pan2026sop_2601_03044,wang2026learning_2605_00416}.

\subsection{From Fixed Data Engines to Autonomous Preparation Agents}

Most current engines automate a controller designed by people: which buffer to
sample, which failures trigger intervention, which generator to call, and when to
retrain are fixed in advance.  The L4 problem in
Figure~\ref{fig:autonomy-ladder} is harder: an agent must select preparation
operators and information-gathering experiments under delayed, noisy utility.
For example, after detecting weak performance on deformable objects, it must
choose among collecting demonstrations, preserving failed prefixes, requesting
corrections, generating scenes, changing a mixture, calibrating a verifier, or
first running a diagnostic rollout.  These actions have different cost, delay,
identifiability, and safety.

This can be studied without granting unrestricted real-robot autonomy.  A
benchmark should expose a versioned raw pool, typed operator library, fixed
learners, simulator and limited hidden real audits, a cost model, and planted
defects or coverage gaps.  At each round an agent receives only the evidence it
purchased, emits an executable preparation plan and predicted gain, and may
promote, abstain, stop, or roll back.  Comparisons should include a fixed expert
pipeline, random and bandit allocation, an oracle with defect labels, and an
LLM planner without learner/rollout feedback.  Metrics are cumulative utility
regret, evidence and robot cost, safety violations, worst-group retention,
calibration of predicted gains, unnecessary rollbacks, and exact replayability
from lineage.

Four capabilities are falsifiable milestones.  \emph{Operator competence}
requires valid, reversible transformations.  \emph{Experimental competence}
requires isolating a data decision with matched controls rather than changing
the learner simultaneously.  \emph{Transfer competence} requires the chosen
pipeline to help a held-out learner or deployment shift.  \emph{Governance
competence} requires respecting authority boundaries, immutable anchors, and
safety budgets even when a proxy predicts a gain.  An LLM that writes a fluent
pipeline but cannot meet these criteria is an interface to preparation, not an
autonomous preparation agent.

\subsection{Cost--Benefit and Environmental Accounting}

Preparation papers rarely include full costs, making an expensive selector appear
data-efficient.  We need standardized accounts for human review, model inference,
subset-training runs, simulation, storage/decode, network movement, target-robot
hours, interventions, and discarded collection.  Benefits should be expressed as
utility reached at a total cost and as amortization across downstream models.

This supports principled cascades.  For example, a deterministic gate may remove
known corruption, a small model may route obvious semantic cases, and influence
or rollout evidence may be purchased only for uncertain high-value groups.  The
optimal cascade is context-dependent; benchmark tracks should therefore impose
several resource profiles rather than one hidden budget.

\subsection{Bias, Consent, Licensing, and Auditability}

Human and community data introduce privacy, consent, worker, and geographic
representation questions; web video introduces licensing and identity risk;
foundation-model filters can systematically reject unfamiliar homes, tools, or
bodies; autonomous collection can record bystanders or unsafe attempts.  These
are preparation decisions because filters, retention policies, provenance, and
access control determine what enters training.

Datasheets should expose source license and consent, sensitive fields, demographic
or environment proxies where lawful, filter rejection rates by subgroup, and
known blind spots.  Provenance must survive synthesis: a generated derivative
should retain source and model lineage.  Utility optimization should include
worst-group and safety constraints so rare environments are not removed merely
because they are harder for a pretrained judge.

\subsection{EvoDataPrep: A Unifying System Hypothesis}

We propose \evodataprep as a research framework, not a claim of a completed
system.  It comprises nine modules: (1) immutable raw store and lineage graph;
(2) typed operator library; (3) hierarchical data units; (4) multidimensional
quality estimators; (5) learner-conditional utility model; (6) role and mixture
router; (7) generator/augmenter; (8) layered independent verifiers; and (9) an
active controller that allocates collection, simulation, labeling, and training
budget.

Given a context and budget, the controller selects a pipeline and routing policy
from composable operators, observes fixed-budget learner and rollout outcomes,
updates its utility model, and promotes a version only if it improves a
multi-objective safety-constrained criterion.  Every decision is reversible
because the raw data and transformation graph remain immutable.  A practical
first study can restrict the search space to segmentation, three curation signals,
mixture weights, and a synthetic-data gate on an open VLA learner.  It should use
multiple contamination and target shifts, compare Bayesian optimization or
bandit allocation with hand pipelines, and measure utility regret plus total cost.

The stronger long-term hypothesis is that data preparation can be learned as a
policy while remaining auditable---the transition from the L3 engine to the L4
agent.  It is falsified if learned pipelines fail to transfer beyond their
calibration learner, consume more total resources than fixed baselines, cannot
attribute their gains to a data decision, or amplify errors across rounds.
Stating these failure criteria is important: ``automation'' alone is not an
academic contribution unless it improves the correctness--coverage--cost frontier
under controlled evidence and can safely decline authority it has not earned.
